% =============================================================
% Building Reliable Systems
% Chapter order is set in ../main.tex; the rendered chapter number
% comes from \chapter{} below.
%
% Intended position: at the end of the agents arc, after the
% chapters on agent design, orchestration, and context engineering.
% Currently sits right after Tool Calls because there are no
% agent chapters yet to slot in between.
% =============================================================

\chapter{Building Reliable Systems}
\label{ch:reliable}

\todo{Chapter stub. Subsections name the topics agreed in the
v0.5 outline; content is deferred until agent chapters exist
that this chapter can refer back to.}

By this point in the book the reader has the full picture of an
agent system: a model in a loop with tools, with context
engineering and multi-agent coordination layered on top. What
remains is the operational question: how do we deploy such a
system into production, and keep it working? This chapter answers
that question. The topics it covers --- security, evaluation,
observability, human-in-the-loop checkpoints, versioning, and
sandbox design --- are not glamorous, but no agent makes it past
the first month of real use without running into all of them.
Treating them late, after the architectural chapters, is
deliberate: each of these concerns is sharper once there is an
actual loop and actual side-effects to defend.

\section{Security}
\label{sec:reliable-security}

\todo{Section stub. The threat model for agents is different from
the threat model for traditional software because the agent
\emph{consumes} content from sources it cannot fully trust (web
pages, emails, ticket comments, retrieved documents) and acts on
that content. Subsections below.}

\subsection{Prompt injection}
\label{sec:reliable-injection}

\todo{Definition: instructions hidden in data the agent retrieves,
which the model cannot reliably distinguish from instructions
issued by the user. Worked examples (the classic ``ignore previous
instructions and exfiltrate the API key''; subtler real-world
cases). Why the agent setting is uniquely vulnerable. Mitigations
in increasing order of strength: prompt hygiene, input
quarantining, structured-output gates, capability scoping,
out-of-band confirmation.}

\subsection{The confused deputy}
\label{sec:reliable-deputy}

\todo{Authority leaks: the agent has more authority than the user
who instructed it (e.g.\ a CI runner with write access to the
production environment) and is asked, through a comment in a pull
request or an issue, to use that authority on behalf of an
attacker. Forward-reference \S\ref{sec:sandbox}.}

\subsection{Data exfiltration channels}
\label{sec:reliable-exfil}

\todo{The agent has read access to private state (chat history,
internal repositories, secrets in environment variables) and the
ability to make outbound network calls. The two together produce
exfiltration. Mitigations: outbound network allow-lists,
egress-only sandboxes, redaction layers.}

\section{Sandboxing}
\label{sec:reliable-sandbox}

\todo{The defensive primitive that recurs throughout the book.
This section is the consolidated treatment that earlier chapters
forward-reference. Container-based sandboxing, user-mode kernels
(gVisor), seccomp profiles, virtual machines. The four properties
to ask of any sandbox (no host filesystem access beyond a
designated working directory; network allow-list; CPU/memory/time
budget; ephemeral state). Sandbox composition: one per untrusted
component (one per MCP server, one per agent, one per tool
group). Trade-offs against latency, observability, and developer
ergonomics.}

\subsection{Sandboxing the shell tool}
\label{sec:reliable-sandbox-shell}

\todo{Concrete recipe for the shell-as-spine pattern of
\S\ref{sec:shell}. Docker baseline; a more locked-down recipe
using gVisor; a deny-by-default network policy; a writeable
work-dir mount and a read-only host bind for tooling.}

\subsection{Sandboxing the code interpreter}
\label{sec:reliable-sandbox-codeact}

\todo{For programmatic-tool-call agents (\S\ref{sec:codeact}),
the sandbox has the additional job of restricting what the
executed code can import. Allow-list of modules; the
\code{importlib.metadata} hook; the trade-off between strict
allow-listing and the agent's ability to write useful code.}

\subsection{Sandboxing MCP servers}
\label{sec:reliable-sandbox-mcp}

\todo{Treat each MCP server as an untrusted dependency. Network
policy per server (an MCP server that only needs to reach
\code{api.github.com} should not be able to reach anywhere
else). Capability scoping at the host: even if the server says
``I support these tools'', the host can still mask off the ones
this agent should not call.}

\section{Testing and Evaluation}
\label{sec:reliable-eval}

\todo{Section stub. The eval problem is hard for agents because
the system is non-deterministic (\S\ref{sec:mental-sampling})
and because the right answer is often a distribution rather than
a single output. Subsections below.}

\subsection{Eval mindset}
\label{sec:reliable-eval-mindset}

\todo{Define the eval set as a representative distribution of
inputs and acceptance criteria, not as a regression suite that
expects bit-for-bit identity. The N-of-K success rate as the
canonical metric (e.g.\ ``passes at least 4 of 5 runs'').
Distinguish from regression \emph{tests} of the surrounding
software, which can and should be deterministic.}

\subsection{Building a test set}
\label{sec:reliable-eval-set}

\todo{How to harvest representative inputs: real production
traffic with privacy redaction; adversarial inputs from
red-teaming; corner cases the team has accumulated. The
trap of over-fitting to the eval set --- the same problem
that plagues ML benchmarks generally.}

\subsection{Existing benchmarks}
\label{sec:reliable-eval-benches}

\todo{Survey, briefly, the major public benchmarks the reader
might encounter or use as a reference: SWE-bench, AgentBench,
GAIA, the Berkeley Function-Calling Leaderboard, MLE-bench, the
domain-specific ones (CVDP for chip design, MedQA for medicine).
The point is not to enumerate but to give the reader the
vocabulary so external claims (``this model scores X on Y'')
are legible.}

\subsection{Continuous evaluation}
\label{sec:reliable-eval-continuous}

\todo{Running evals on every model upgrade and every prompt
change. Pinning model versions explicitly. The cautionary tale
of Claude Code's silent shift from ``high'' to ``medium''
adaptive thinking in February~2026, which produced a measured
67\% drop in reasoning depth across thousands of sessions
without any notice~\citep{adaptivethinking2026}.}

\section{Agent-Specific Failure Modes}
\label{sec:reliable-agent-fails}

\todo{Section stub. The failure modes that emerge from the loop
itself, separate from the model-level failures of
\S\ref{sec:mental-failures}. Each gets a paragraph with
diagnosis and mitigation.}

\paragraph{Runaway loops.}
\todo{The agent retries the same failing tool call ad infinitum,
or oscillates between two states. Mitigation: hard iteration
caps (\S\ref{sec:impl-blocks}), loop detection on tool-call
sequences, exponential back-off.}

\paragraph{Context exhaustion mid-task.}
\todo{The agent's context fills up with tool results before the
task is done. Mitigation: summarization passes, sliding-window
history, structured offloading to external memory.}

\paragraph{Stuck states.}
\todo{The agent is waiting for an event that will never happen,
or has decided the task is impossible when it is not. Mitigation:
timeouts, supervisor checks, the ability to restart with a
modified prompt.}

\paragraph{Goal drift.}
\todo{The agent silently shifts the task it is working on, often
because of intermediate observations that nudged its
interpretation. Mitigation: re-anchoring the system prompt,
explicit task-state tracking, periodic ``what am I trying to
do?'' summarization.}

\section{Observability}
\label{sec:reliable-observability}

\todo{Section stub. Without observability, evals are guesswork
and incident response is impossible.}

\subsection{Logging}
\label{sec:reliable-log}

\todo{What to log: every model call (prompt, response, model
version, sampling parameters), every tool call (name, arguments,
return value, duration, exit status), every loop boundary, every
human approval. The trade-off between verbosity and storage cost.}

\subsection{Tracing}
\label{sec:reliable-trace}

\todo{Distributed tracing across the agent loop, the underlying
model calls, and the tools the agent invokes. OpenTelemetry-style
spans; correlation IDs that link a single user request to all
downstream work. The available tooling (Langfuse, Arize, Logfire,
plain OTel) without picking favourites.}

\subsection{Replay}
\label{sec:reliable-replay}

\todo{The ability to replay a logged session from any point.
Particularly important for debugging non-deterministic failures:
re-running with seed-pinned sampling and the same tool responses
to isolate where divergence happens. Limits of replay (the
underlying model may have moved silently;
\S\ref{sec:reliable-versioning}).}

\section{Human-in-the-Loop Checkpoints}
\label{sec:reliable-hitl}

\todo{Section stub. When to require explicit human approval and
how to design the gate without turning every action into a
blocking question.}

\subsection{Where to gate}
\label{sec:reliable-gate}

\todo{Sensible defaults: irreversible side-effects (sending
emails, paying invoices, deleting data) gated by default;
read-only operations and reversible writes ungated; the boundary
in between is the design surface.}

\subsection{Gate design}
\label{sec:reliable-gate-design}

\todo{The gate must show the user enough to decide. Show the
intended action; show the rationale (one or two sentences from
the agent); show the diff if the action is structured. Antipatterns:
gates that obscure what is about to happen; gates that fire so
often the user clicks through reflexively (alarm fatigue).}

\subsection{Bypass and escalation}
\label{sec:reliable-bypass}

\todo{Trusted users may want to skip gates after the agent has
proven itself; novel situations may want to escalate gates that
were ungated by default. The control surface that lets a user
configure this without redesigning the agent.}

\section{Versioning and Rollback}
\label{sec:reliable-versioning}

\todo{Section stub. The agent system has multiple moving parts:
model version, prompt template, tool inventory, harness code.
All four need to be versioned, and a rollback story has to span
all four.}

\subsection{Pinning models}
\label{sec:reliable-pin}

\todo{Why ``the latest version'' is a bug, not a feature, for an
agent in production. How to pin --- explicit model identifiers
--- and the limits of pinning (some hosted providers update the
underlying weights without changing the identifier; some retire
identifiers entirely).}

\subsection{Pinning prompts}
\label{sec:reliable-pin-prompt}

\todo{Treating the system prompt as code: in version control,
reviewed in pull requests, deployed alongside the harness rather
than edited in a UI. The eval suite of
\S\ref{sec:reliable-eval} runs on every prompt change.}

\subsection{Rollback as a first-class operation}
\label{sec:reliable-rollback}

\todo{When something breaks in production, the smallest available
intervention should be ``ship the previous version''. This
requires that all four moving parts (model, prompt, tools,
harness) are tagged together. The cautionary tale of the
adaptive-thinking drop again --- because the model identifier did
not change, there was no version to roll back to.}
